\documentclass[10pt,a4paper]{amsart}
\usepackage{amsmath,amssymb,amsthm,graphicx}
% Retained minimal preamble for the Stacks Project chapter more-algebra.tex (More on Algebra).
% The chapter file pinned-original-source.tex carries only \input{preamble} on its first line, so no
% preamble text lives inside the chapter itself.  The Stacks Project uses one single shared preamble for
% every chapter; the snapshot stages it identically as sources/stacks-chow/preamble.tex and
% sources/stacks-etale-cohomology/preamble.tex.  The block below is copied mechanically, line for line,
% from sources/stacks-chow/preamble.tex lines 14-18 (the Xy-pic setup) and 154-217 (the native theorem
% declarations together with the mathematical macros).  Nothing was added, removed or redefined; the
% wrapper class and the amsmath/amssymb/amsthm/graphicx base packages are supplied by the assembler.
% Chapter: more-algebra.tex -- More on Algebra.
% For commutative diagrams we use Xy-pic
\usepackage[all]{xy}

% We use 2cell for 2-commutative diagrams.
\xyoption{2cell}

% Theorem environments.
%
\theoremstyle{plain}
\newtheorem{theorem}[subsection]{Theorem}
\newtheorem{proposition}[subsection]{Proposition}
\newtheorem{lemma}[subsection]{Lemma}

\theoremstyle{definition}
\newtheorem{definition}[subsection]{Definition}
\newtheorem{example}[subsection]{Example}
\newtheorem{exercise}[subsection]{Exercise}
\newtheorem{situation}[subsection]{Situation}

\theoremstyle{remark}
\newtheorem{remark}[subsection]{Remark}
\newtheorem{remarks}[subsection]{Remarks}

\numberwithin{equation}{subsection}

% Macros
%
\def\lim{\mathop{\mathrm{lim}}\nolimits}
\def\colim{\mathop{\mathrm{colim}}\nolimits}
\def\Spec{\mathop{\mathrm{Spec}}}
\def\Hom{\mathop{\mathrm{Hom}}\nolimits}
\def\Ext{\mathop{\mathrm{Ext}}\nolimits}
\def\SheafHom{\mathop{\mathcal{H}\!\mathit{om}}\nolimits}
\def\SheafExt{\mathop{\mathcal{E}\!\mathit{xt}}\nolimits}
\def\Sch{\mathit{Sch}}
\def\Mor{\mathop{\mathrm{Mor}}\nolimits}
\def\Ob{\mathop{\mathrm{Ob}}\nolimits}
\def\Sh{\mathop{\mathit{Sh}}\nolimits}
\def\NL{\mathop{N\!L}\nolimits}
\def\CH{\mathop{\mathrm{CH}}\nolimits}
\def\proetale{{pro\text{-}\acute{e}tale}}
\def\etale{{\acute{e}tale}}
\def\QCoh{\mathit{QCoh}}
\def\Ker{\mathop{\mathrm{Ker}}}
\def\Im{\mathop{\mathrm{Im}}}
\def\Coker{\mathop{\mathrm{Coker}}}
\def\Coim{\mathop{\mathrm{Coim}}}

% Boxtimes
%
\DeclareMathSymbol{\boxtimes}{\mathbin}{AMSa}{"02}

%
% Macros for moduli stacks/spaces
%
\def\QCohstack{\mathcal{QC}\!\mathit{oh}}
\def\Cohstack{\mathcal{C}\!\mathit{oh}}
\def\Spacesstack{\mathcal{S}\!\mathit{paces}}
\def\Quotfunctor{\mathrm{Quot}}
\def\Hilbfunctor{\mathrm{Hilb}}
\def\Curvesstack{\mathcal{C}\!\mathit{urves}}
\def\Polarizedstack{\mathcal{P}\!\mathit{olarized}}
\def\Complexesstack{\mathcal{C}\!\mathit{omplexes}}
% \Pic is the operator that assigns to X its picard group, usage \Pic(X)
% \Picardstack_{X/B} denotes the Picard stack of X over B
% \Picardfunctor_{X/B} denotes the Picard functor of X over B
\def\Pic{\mathop{\mathrm{Pic}}\nolimits}
\def\Picardstack{\mathcal{P}\!\mathit{ic}}
\def\Picardfunctor{\mathrm{Pic}}
\def\Deformationcategory{\mathcal{D}\!\mathit{ef}}
\usepackage{cleveref}
\crefname{originallemma}{original lemma}{original lemmas}
\Crefname{originallemma}{Original lemma}{Original lemmas}
\makeatletter
\newlabel{more-algebra-theorem-regular-fs}{{07EL}{source}{}{sourcepointer-1}{}}
\newlabel{more-algebra-theorem-regular-fs@cref}{{[originallemma][0][]07EL}{[1][0][]source}{}{}{}}
\makeatother
\crefname{originallemma}{original lemma}{original lemmas}
\Crefname{originallemma}{Original lemma}{Original lemmas}
\makeatletter
\newlabel{lemma-fs-implies-regular}{{07EI}{source}{}{sourcepointer-2}{}}
\newlabel{lemma-fs-implies-regular@cref}{{[originallemma][0][]07EI}{[1][0][]source}{}{}{}}
\makeatother
\crefname{originallemma}{original lemma}{original lemmas}
\Crefname{originallemma}{Original lemma}{Original lemmas}
\makeatletter
\newlabel{lemma-regular-implies-fs}{{07EJ}{source}{}{sourcepointer-3}{}}
\newlabel{lemma-regular-implies-fs@cref}{{[originallemma][0][]07EJ}{[1][0][]source}{}{}{}}
\makeatother
\crefname{originallemma}{original lemma}{original lemmas}
\Crefname{originallemma}{Original lemma}{Original lemmas}
\makeatletter
\newlabel{proposition-characterization-geometrically-regular}{{07E5}{source}{}{sourcepointer-4}{}}
\newlabel{proposition-characterization-geometrically-regular@cref}{{[originallemma][0][]07E5}{[1][0][]source}{}{}{}}
\makeatother
\crefname{originallemma}{original lemma}{original lemmas}
\Crefname{originallemma}{Original lemma}{Original lemmas}
\makeatletter
\newlabel{lemma-base-change-fs}{{07EG}{source}{}{sourcepointer-5}{}}
\newlabel{lemma-base-change-fs@cref}{{[originallemma][0][]07EG}{[1][0][]source}{}{}{}}
\makeatother
\crefname{originallemma}{original lemma}{original lemmas}
\Crefname{originallemma}{Original lemma}{Original lemmas}
\makeatletter
\newlabel{definition-formally-smooth}{{07EB}{source}{}{sourcepointer-6}{}}
\newlabel{definition-formally-smooth@cref}{{[originallemma][0][]07EB}{[1][0][]source}{}{}{}}
\makeatother
\crefname{originallemma}{original lemma}{original lemmas}
\Crefname{originallemma}{Original lemma}{Original lemmas}
\makeatletter
\newlabel{algebra-lemma-universal-omega}{{00RO}{source}{}{sourcepointer-7}{}}
\newlabel{algebra-lemma-universal-omega@cref}{{[originallemma][0][]00RO}{[1][0][]source}{}{}{}}
\makeatother
\title{A regular ring of finite type over a field is formally smooth after inverting one element}
\author{The Stacks Project Authors}
\date{stacks-project:more-algebra.tex:local-snapshot-20261012; GNU FDL 1.2 or later}
\begin{document}
\maketitle
\setcounter{subsection}{1}
\noindent\textit{Editorial note.} All delivered mathematical text is the author's own, copied from the pinned Stacks Project chapter file more-algebra.tex; no statement or proof step was rewritten, reordered or removed.\par\medskip
\noindent\textit{Editorial note.} The chapter file pinned-original-source.tex carries only \textbackslash{}input\{preamble\} on its first line, so no preamble text lives inside the chapter. The retained file sources/stacks-more-algebra/retained-preamble.tex is a mechanical, line-for-line copy of the project's single shared preamble as staged in the snapshot (Xy-pic setup plus theorem declarations and mathematical macros). Nothing was added, removed or redefined, and the excerpt serves only as preamble, never as a span coordinate.\par\medskip
\noindent\textit{Editorial note.} Both counted spans start at the \textbackslash{}begin line of their environment, so the statement and the author's proof are delivered whole; surrounding prose of the section is deliberately outside the counted spans and is carried as standard prerequisites instead.\par\medskip
\noindent\textit{Editorial note.} The citation fields (source\_id, source\_version, authors, source\_url, license.name, source\_title) are identical to those of every other card this lane has produced from this chapter, so the batch selector treats them as one source.\par\medskip
\noindent\textit{Allowed standard prerequisites (editorial).}\begin{itemize}
\item Formal smoothness and the Jacobian criterion for a finite type algebra over a field. used in both cases, proof lines 9894-9951. Supplies the criterion that has to be satisfied after inverting the element.
\item Regularity of a finite type algebra over a field and its consequences for the module of differentials. used at proof lines 9894-9920. Provides the algebraic input that makes the differential criterion applicable.
\item p-basis theory in characteristic p, developed in the same chapter. the statement counted by the sibling card S1-algebra-p-basis at 11362-11415. Supplies the characteristic p case of the proof.
\end{itemize}\medskip
\ifdefined\hypertarget\hypertarget{sourcepointer-1}{}\fi
\noindent\textit{Source reference.} Original lemma 07EL: focused-seeds/AGP/sources/stacks-snapshot-20261011/tags.txt, source lines 9451--9451. The lemma's own canonical tag, quoted so that readers can resolve the label inside the delivered text; it points at the counted statement itself and carries no proof content.\par\medskip
\ifdefined\hypertarget\hypertarget{sourcepointer-2}{}\fi
\noindent\textit{Source reference.} Original lemma 07EI: focused-seeds/AGP/sources/stacks-snapshot-20261011/tags.txt, source lines 9448--9448. Navigation pointer to a statement the counted proof refers to but does not contain; nothing load-bearing is delegated to it.\par\medskip
\ifdefined\hypertarget\hypertarget{sourcepointer-3}{}\fi
\noindent\textit{Source reference.} Original lemma 07EJ: focused-seeds/AGP/sources/stacks-snapshot-20261011/tags.txt, source lines 9449--9449. Navigation pointer to a statement the counted proof refers to but does not contain; nothing load-bearing is delegated to it.\par\medskip
\ifdefined\hypertarget\hypertarget{sourcepointer-4}{}\fi
\noindent\textit{Source reference.} Original lemma 07E5: focused-seeds/AGP/sources/stacks-snapshot-20261011/tags.txt, source lines 9435--9435. Navigation pointer to a statement the counted proof refers to but does not contain; nothing load-bearing is delegated to it.\par\medskip
\ifdefined\hypertarget\hypertarget{sourcepointer-5}{}\fi
\noindent\textit{Source reference.} Original lemma 07EG: focused-seeds/AGP/sources/stacks-snapshot-20261011/tags.txt, source lines 9446--9446. Navigation pointer to a statement the counted proof refers to but does not contain; nothing load-bearing is delegated to it.\par\medskip
\ifdefined\hypertarget\hypertarget{sourcepointer-6}{}\fi
\noindent\textit{Source reference.} Original lemma 07EB: focused-seeds/AGP/sources/stacks-snapshot-20261011/tags.txt, source lines 9441--9441. Navigation pointer to a statement the counted proof refers to but does not contain; nothing load-bearing is delegated to it.\par\medskip
\ifdefined\hypertarget\hypertarget{sourcepointer-7}{}\fi
\noindent\textit{Source reference.} Original lemma 00RO: focused-seeds/AGP/sources/stacks-snapshot-20261011/tags.txt, source lines 1010--1010. Navigation pointer to a statement the counted proof refers to but does not contain; nothing load-bearing is delegated to it.\par\medskip
\begin{theorem}
\label{theorem-regular-fs}
Let $k$ be a field. Let $(A, \mathfrak m, K)$ be a Noetherian local
$k$-algebra. If the characteristic of $k$ is zero then the following
are equivalent
\begin{enumerate}
\item $A$ is a regular local ring, and
\item $k \to A$ is formally smooth in the $\mathfrak m$-adic topology.
\end{enumerate}
If the characteristic of $k$ is $p > 0$ then the following are equivalent
\begin{enumerate}
\item $A$ is geometrically regular over $k$,
\item $k \to A$ is formally smooth in the $\mathfrak m$-adic topology.
\item for all $k \subset k' \subset k^{1/p}$
finite over $k$ the ring $A \otimes_k k'$ is regular,
\item $A$ is regular and the canonical map
$H_1(L_{K/k}) \to \mathfrak m/\mathfrak m^2$ is injective, and
\item $A$ is regular and the map
$\Omega_{k/\mathbf{F}_p} \otimes_k K \to \Omega_{A/\mathbf{F}_p} \otimes_A K$
is injective.
\end{enumerate}
\end{theorem}

\begin{proof}
If the characteristic of $k$ is zero, then the equivalence of (1) and (2)
follows from
Lemmas \ref{lemma-fs-implies-regular} and \ref{lemma-regular-implies-fs}.

\medskip\noindent
If the characteristic of $k$ is $p > 0$, then it follows from
Proposition \ref{proposition-characterization-geometrically-regular}
that (1), (3), (4), and (5) are equivalent. Assume (2) holds.
By Lemma \ref{lemma-base-change-fs} we see that
$k' \to A' = A \otimes_k k'$ is formally smooth for the
$\mathfrak m' = \mathfrak mA'$-adic topology. Hence if $k \subset k'$ is
finite purely inseparable, then $A'$ is a regular local ring by
Lemma \ref{lemma-fs-implies-regular}. Thus we see that (1) holds.

\medskip\noindent
Finally, we will prove that (5) implies (2). Choose a solid diagram
$$
\xymatrix{
A \ar[r]_{\bar\psi} \ar@{-->}[rd] & B/J \\
k \ar[u]^i \ar[r]^\varphi & B \ar[u]_\pi
}
$$
as in Definition \ref{definition-formally-smooth}. As $J^2 = 0$ we see
that $J$ has a canonical $B/J$ module structure and via $\bar\psi$ an
$A$-module structure. As $\bar\psi$ is continuous for the
$\mathfrak m$-adic topology we see that $\mathfrak m^nJ = 0$ for some $n$.
Hence we can filter $J$ by $B/J$-submodules
$0 \subset J_1 \subset J_2 \subset \ldots \subset J_n = J$
such that each quotient $J_{t + 1}/J_t$ is annihilated by $\mathfrak m$.
Considering the sequence of ring maps
$B \to B/J_1 \to B/J_2 \to \ldots \to B/J$
we see that it suffices to prove the existence of the dotted arrow when
$J$ is annihilated by $\mathfrak m$, i.e., when $J$ is a $K$-vector space.

\medskip\noindent
Assume given a diagram as above such that $J$ is annihilated by $\mathfrak m$.
By Lemma \ref{lemma-regular-implies-fs} we see that $\mathbf{F}_p \to A$ is
formally smooth in the $\mathfrak m$-adic topology. Hence we can find a ring
map $\psi : A \to B$ such that $\pi \circ \psi = \bar \psi$. Then
$\psi \circ i, \varphi : k \to B$ are two maps whose compositions with $\pi$
are equal. Hence $D = \psi \circ i - \varphi : k \to J$ is a derivation.
By Algebra, Lemma \ref{algebra-lemma-universal-omega} we can write
$D = \xi \circ \text{d}$ for some $k$-linear map
$\xi : \Omega_{k/\mathbf{F}_p} \to J$. Using the $K$-vector space structure
on $J$ we extend $\xi$ to a $K$-linear map
$\xi' : \Omega_{k/\mathbf{F}_p} \otimes_k K \to J$.
Using (5) we can find a $K$-linear map
$\xi'' : \Omega_{A/\mathbf{F}_p} \otimes_A K$ whose restriction to
$\Omega_{k/\mathbf{F}_p} \otimes_k K$ is $\xi'$. Write
$$
D' : A \xrightarrow{\text{d}} \Omega_{A/\mathbf{F}_p}
\to \Omega_{A/\mathbf{F}_p} \otimes_A K \xrightarrow{\xi''} J.
$$
Finally, set $\psi' = \psi - D' : A \to B$. The reader verifies that $\psi'$
is a ring map such that $\pi \circ \psi' = \bar \psi$ and such that
$\psi' \circ i = \varphi$ as desired.
\end{proof}

\end{document}
